% संपूर्ण मराठी ग्रंथातील प्रकरण 51: ज्ञानविषयक तर्कशास्त्रे
% हा संपादनयोग्य स्रोतखंड आहे. संकलनासाठी संपूर्ण स्रोतसंग्रहातील
% अक्षररूपे, पूर्वभाग, चिन्हव्याख्या आणि आकृती-स्रोत आवश्यक आहेत.
% मूळ: Open Logic Project; मजकूर आणि रूपांतर CC BY 4.0.
% यंत्रानुवाद, दुरुस्ती आणि तपासणी: OpenAI Codex —
% GPT-5.6 Sol आणि GPT-6 Sol, दोन्ही Ultra effort.
% दुरुस्ती-पुनरावलोकन: GPT-6.1 Sol, Ultra effort.
\chapter{ज्ञानविषयक तर्कशास्त्रे}\label{aml:el:chap}
\begin{editorial}
  या प्रकरणात ज्ञानविषयक तर्कशास्त्रांचा अधिसिद्धांत
  मांडला आहे. याची रचना अल्डो अँटोनेली यांच्या
  अभिजात मूलभूत मोडल तर्कशास्त्रावरील नोंदींसारखी
  आहे; मात्र द्विअनुकरण आणि गतिशील ज्ञानविषयक
  तर्कशास्त्रांवरील साहित्य जोडण्यासाठी ऑड्री यॅप
  यांनी त्याचे पुनर्लेखन केले आहे.
\end{editorial}
\section{प्रस्तावना}\label{aml:el:int:sec}
मोडल तर्कशास्त्र जसे \emph{मोडल विधानांचा} आणि
त्यांमधील तार्किक निष्पन्नतेच्या संबंधांचा विचार
करते, तसे ज्ञानविषयक तर्कशास्त्र \emph{ज्ञानविषयक
विधानांचा} आणि त्यांमधील निष्पन्नता संबंधांचा विचार
करते. मोडल संकारकांचा अर्थ संभाव्यता आणि
अनिवार्यता असा लावण्याऐवजी, एकस्थानी संयोजकांना
ज्ञानविषयक किंवा विश्वासविषयक अर्थ दिला जातो;
त्यायोगे ज्ञान आणि विश्वासाचे प्रतिमानीकरण करता
येते. उदाहरणार्थ, पुढील विधाने व्यक्त करायची
असू शकतात:
\begin{enumerate}
\item कॅलगरी हे अल्बर्टात आहे, हे रिचर्डला माहीत आहे.
\item सोफ्यावर कुत्रा असू शकतो, असे ऑड्रीला वाटते.
\item मंगळवारी तिचा वर्ग असतो, हे ऑड्रीला माहीत
  आहे, हे रिचर्डला माहीत आहे.
\item एका वर्षात १२ महिने असतात, हे सर्वांना माहीत आहे.
\end{enumerate}
आधुनिक ज्ञानविषयक तर्कशास्त्राचा उगम अनेकदा
याक्को हिंटिक्का यांच्या १९६२ मधील \emph{Knowledge
and Belief} या ग्रंथाशी जोडला जातो. तो ग्रंथ
लिहिला गेला तेव्हा संभाव्य जगांचे अर्थनिर्धारण
तर्कशास्त्रात वाढत्या प्रमाणात वापरले जात होते.
वास्तविक, ज्ञानविषयक तर्कशास्त्रे इतर मोडल
तर्कशास्त्रांसारखीच बहुतेक अर्थनिर्धारणाची साधने
वापरतात; मात्र त्यांचा अर्थ वेगळा लावतात. मुख्य
बदल म्हणजे \emph{प्राप्यता संबंध} कशाचे प्रतिनिधित्व
करतो, हे. ज्ञानविषयक तर्कशास्त्रांत तो काही प्रकारची
\emph{ज्ञानविषयक संभाव्यता} दर्शवतो. आपण कोणत्या
ज्ञानविषयक संकल्पनेचे प्रतिमानीकरण करतो त्यावर
प्राप्यता संबंधावर कोणत्या अटी घालायच्या हे अवलंबून
असेल, हे पुढे पाहू. अनुरूपता सिद्धांताला ज्ञानविषयक
अर्थ दिल्यावर काय घडते तेही पाहू. वरील उदाहरणांत
रिचर्ड आणि ऑड्री हे दोन कर्ते आहेत, आणि प्रत्येकीला
काय माहीत आहे यांतील संबंध येतो. आपण विचारात
घेणार असलेल्या ज्ञानविषयक तर्कशास्त्रांत अनेक कर्ते
असतील, म्हणून अशा गोष्टी व्यक्त करता येतील.
याउलट, एका कर्त्याचे ज्ञानविषयक तर्कशास्त्र फक्त
एकाच व्यक्तीला काय माहीत आहे किंवा काय वाटते
याबद्दल बोलेल.
\section{ज्ञानविषयक तर्कशास्त्राची भाषा}\label{aml:el:lan:sec}
\begin{defn}
$G$ हा कर्तृचिन्हांचा संच असू द्या. अनेक कर्त्यांच्या
ज्ञानविषयक तर्कशास्त्राच्या मूलभूत भाषेत पुढील गोष्टी
असतात:
\begin{enumerate}
  \item असत्यता दर्शवणारा विधानीय स्थिरांक~$\lfalse$.
  \item सत्यता दर्शवणारा विधानीय स्थिरांक~$\ltrue$.
  \item विधानीय चलांचा गणनीय अनंत संच: $\Obj
    p_0$, $\Obj p_1$, $\Obj p_2$, \dots
  \item विधानीय संयोजके: \startycommalist
  \ycomma $\lnot$ (नकार)
  \ycomma $\land$ (संयोग)
  \ycomma $\lor$ (विकल्पयोग)
  \ycomma $\lif$ (शर्त)
  \ycomma $\liff$ (द्विशर्त).
  \item $a \in G$ असताना ज्ञान-संकारक $\Knows_a$.
\end{enumerate}
\end{defn}
आपल्या व्यवस्थेत आपल्याला एकाच कर्त्याच्या ज्ञानाचा
विचार करायचा असेल, तर संच~$G$ आणि स्वतंत्र कर्ते
यांचा निर्देश वगळता येतो. त्या वेळी आपल्याकडे फक्त
मूलभूत संकारक~$\Knows$ असतो.
\begin{defn}
ज्ञानविषयक भाषेतील \emph{सूत्रे} ची पुढीलप्रमाणे
  आगमनाने व्याख्या करतो:
\begin{enumerate}
\item $\lfalse$ हे अण्विक सूत्र आहे.
\item $\ltrue$ हे अण्विक सूत्र आहे.
\item प्रत्येक विधानीय चल~$\Obj p_i$ हे (अण्विक) सूत्र आहे.
\item $\varphi$ हे सूत्र असेल, तर $\lnot \varphi$
  हेही सूत्र आहे.
\item $\varphi$ आणि $\psi$ ही सूत्रे असतील, तर $(\varphi \land
  \psi)$ हे सूत्र आहे.
\item $\varphi$ आणि $\psi$ ही सूत्रे असतील, तर $(\varphi \lor \psi)$
  हे सूत्र आहे.
\item $\varphi$ आणि $\psi$ ही सूत्रे असतील, तर $(\varphi \lif \psi)$
  हे सूत्र आहे.
\item $\varphi$ आणि $\psi$ ही सूत्रे असतील, तर $(\varphi \liff \psi)$
  हे सूत्र आहे.
\item $\varphi$ हे सूत्र असेल आणि $a \in G$ असेल,
  तर $\Knows_a \varphi$ हे सूत्र आहे.
\item याखेरीज दुसरे काहीही सूत्र नाही.
\end{enumerate}
\end{defn}
सूत्र~$\varphi$ मध्ये कोणत्याही $a\in G$ साठी
$\Knows_a$ नसेल, तर त्याला \emph{मोडल-विरहित} म्हणतो.
\begin{defn}
  $\Knows$ हा संकारक वैयक्तिक ज्ञान दर्शवण्यासाठी
  आहे; $\EKnows$ हा ``सर्वांना माहीत आहे'' असा
  वाचला जाणारा संकारक समूहज्ञान दर्शवतो. रिकामेतर
  सीमित $G' \subseteq G$ साठी $\EKnows_{G'} \varphi$ हे
  पुढील सीमित संयोगाचे संक्षिप्त रूप म्हणून परिभाषित करतो:
  \[\bigwedge_{b \in G'} \Knows_b \varphi.\]
  रिकाम्या समूहासाठी हे सत्य असलेला रिकामा संयोग मानतो.
  $G'$ अनंत असेल, तर हा संयोग मूलभूत सीमित सूत्रांच्या भाषेतील
  संक्षेप नाही. अशा समूहासाठी स्वतंत्र $\EKnows_{G'}$
  संकारक जोडल्यास त्याची सत्यतेची अट अशी घ्यायची:
  $\mSat{M}{\EKnows_{G'} \varphi}[w]$ जर आणि फक्त जर
  प्रत्येक $b\in G'$ साठी $\mSat{M}{\Knows_b \varphi}[w]$.
\end{defn}

कर्त्यांच्या~$G$ समूहातील ज्ञानाची अधिक प्रबळ
संकल्पनाही विचारात घेता येते: \emph{सामायिक ज्ञान}.
यात त्या समूहातील कर्त्यांच्या प्रत्येक रिकामेतर सीमित
क्रमासाठी, एकाला माहीत आहे की दुसऱ्याला माहीत आहे की
आणखी एकाला माहीत आहे, अशा प्रकारचे ज्ञान आवश्यक असते;
अशी अट प्रत्येक सीमित लांबीसाठी घेतली जाते.
सामायिक ज्ञानातून समूहज्ञान मिळते, पण उलट दिशा सर्व
प्रतिमानांत खरी नसते. तरीही प्रत्येक समूह व चौकटीत
हा भेद काटेकोर असेलच असे नाही; उदाहरणार्थ, एकाच
कर्त्याचा परावर्ती संक्रमक संबंध असेल, तर त्याचे
वैयक्तिक आणि सामायिक ज्ञान जुळते.
``$\varphi$ हे~$G$ समूहात सामायिक ज्ञान आहे'' हे
दर्शवण्यासाठी आपण $\CKnows_G \varphi$ वापरू. सर्व सीमित
लांबींच्या अटींचा संयोग स्वतः मूलभूत भाषेतील सीमित
सूत्र नाही; म्हणून $\CKnows_G$ हा विस्तारित भाषेतील
स्वतंत्र संकारक मानतो. त्याचे संबंधी अर्थनिर्धारण
\ref{aml:el:trw:sec} मध्ये दिले आहे.













\section{संबंधी प्रतिमाने}\label{aml:el:rel:sec}

ज्ञानविषयक तर्कशास्त्रांची मूलभूत चिन्हार्थमीमांसक
संकल्पना सामान्य मोडल तर्कशास्त्रांतील संकल्पनेसारखीच
आहे. संबंधी प्रतिमानांत अजूनही जगांचा संच आणि कोणती
विधानीय चले कोणत्या जगांत ``सत्य''
मानायची हे ठरवणारे मूल्यांकन असते. आपण फक्त एकाच
कर्त्याचा विचार करत असू, तर नेहमीप्रमाणे एकच
प्राप्यता संबंध असतो. मात्र अनेक कर्त्यांच्या ज्ञानविषयक
तर्कशास्त्रात एकाच प्राप्यता संबंधाऐवजी प्राप्यता
संबंधांचा संच असतो: कर्तृचिन्हांच्या~$G$ संचातील
प्रत्येक~$a$ साठी एक संबंध.

\emph{संबंधी प्रतिमानात} जगांचा संच असतो; प्रत्येक
कर्त्यासाठी असलेल्या द्विस्थानी प्राप्यता संबंधांनी
ही जगे संबंधित असतात. त्यात कोणती विधानीय चले कोणत्या जगांत सत्य आहेत हे ठरवणारे
मूल्यांकनही असते.

\begin{defn}
  अनेक कर्त्यांच्या ज्ञानविषयक भाषेचे \emph{प्रतिमान}
  म्हणजे $\mModel{M} = \tuple{W, R, V}$ असे त्रिकूट;
  येथे
  \begin{enumerate}
  \item $W$ हा ``जगांचा'' रिकामा नसलेला संच आहे,
  \item प्रत्येक $a \in G$ साठी ${R}_a$ हा~$W$ वरील
  द्विस्थानी प्राप्यता संबंध आहे; $R=(R_a)_{a\in G}$
  हा या संबंधांचा कर्त्यांनुसार निर्देशित परिवार आहे, आणि
  \item $V$ हे प्रत्येक विधानीय चल~$p$ ला संभाव्य जगांचा संच~$V(p)\subseteq W$ नेमून
    देणारे फलन आहे.
  \end{enumerate}
  $R_a ww'$ असेल तेव्हा $w'$ हे कर्ता~$a$ साठी~$w$
  पासून \emph{प्राप्य} आहे असे म्हणतो. $w \in V(p)$
  असेल तेव्हा $p$ हे~$w$ येथे \emph{सत्य} आहे असे म्हणतो.
\end{defn}

याची कार्यपद्धती सामान्य मोडल तर्कशास्त्रासारखीच
आहे; फक्त आणखी प्राप्यता संबंध जोडलेले असतात.
दिलेल्या कर्त्याचा प्राप्यता संबंध सामान्यतः त्या
कर्त्याच्या माहितीच्या स्थितीबद्दल काहीतरी दर्शवतो,
असा अर्थ आपण लावू. उदाहरणार्थ, $R_a ww'$ हे
अनेकदा $w'$ हे जग~$w$ येथे~$a$ कडे असलेल्या
माहितीशी सुसंगत आहे, असे व्यक्त करते असे मानतो.
दुसऱ्या शब्दांत, $w$ येथे त्या कर्त्याला $w$ आणि
जग~$w'$ यांमधील फरक ओळखता येत नाही.













\section{एखाद्या जगात सत्य असणे}\label{aml:el:trw:sec}

सामान्य मोडल तर्कशास्त्राप्रमाणेच प्रत्येक ज्ञानविषयक
प्रतिमान त्यात कोणती सूत्रे कोणत्या जगांत सत्य
मानायची हे ठरवते. संबंधी चिन्हार्थमीमांसेच्या मूलभूत
संकल्पनेसाठी ``प्रतिमान~$\mModel{M}$ मध्ये
सूत्र~$\varphi$ हे जग~$w$ येथे सत्य आहे'' हीच संज्ञा
वापरतो. या संबंधाची आगमनाने व्याख्या केली जाते;
मोडल नसलेल्या सर्व संकारकांसाठी ती सामान्य मोडल
प्रकरणातील व्याख्येसारखीच असते.

\begin{defn}\label{aml:el:trw:defn:mmodels}
  प्रतिमान~$\mModel M$ मध्ये \emph{सूत्र~$\varphi$ हे~$w$ येथे
  सत्य असणे}, म्हणजे $\mSat{M}{\varphi}[w]$, पुढीलप्रमाणे
  आगमनाने परिभाषित करतो:
  \begin{enumerate}
  \item \indcase{\varphi}{\lfalse}{$\mSat{M}{\lfalse}[w]$ कधीच नसते}.
  \item \indcase{\varphi}{\ltrue}{$\mSat{M}{\ltrue}[w]$ नेहमी असते}.
  \item $\mSat{M}{p}[w]$ तेव्हा आणि केवळ तेव्हाच, जेव्हा $w \in V(p)$
  \item \indcase{\varphi}{\lnot \psi}{$\mSat{M}{\indfrm}[w]$
    तेव्हा आणि केवळ तेव्हाच, जेव्हा $\mSat/{M}{\psi}[w]$}.
  \item \indcase{\varphi}{(\psi \land \chi)}{$\mSat{M}{\indfrm}[w]$
    तेव्हा आणि केवळ तेव्हाच, जेव्हा $\mSat{M}{\psi}[w]$ आणि
    $\mSat{M}{\chi}[w]$}.
  \item \indcase{\varphi}{(\psi \lor \chi)}{$\mSat{M}{\indfrm}[w]$
    तेव्हा आणि केवळ तेव्हाच, जेव्हा $\mSat{M}{\psi}[w]$ किंवा
    $\mSat{M}{\chi}[w]$} (किंवा दोन्ही).
  \item \indcase{\varphi}{(\psi \lif \chi)}{$\mSat{M}{\indfrm}[w]$
    तेव्हा आणि केवळ तेव्हाच, जेव्हा $\mSat/{M}{\psi}[w]$ किंवा
    $\mSat{M}{\chi}[w]$}.
  \item \indcase{\varphi}{(\psi \liff \chi)}{$\mSat{M}{\indfrm}[w]$
    तेव्हा आणि केवळ तेव्हाच, जेव्हा $\mSat{M}{\psi}[w]$ आणि
    $\mSat{M}{\chi}[w]$ ही दोन्ही असतात, किंवा
    $\mSat{M}{\psi}[w]$ आणि $\mSat{M}{\chi}[w]$ यांपैकी
    कोणतेही नसते}.
  \item\label{aml:el:trw:defn:sub:mmodels-box}
    \indcase{\varphi}{\Knows_a \psi}{$\mSat{M}{\indfrm}[w]$ तेव्हा आणि
    केवळ तेव्हाच, जेव्हा $R_a ww'$ असलेल्या प्रत्येक $w' \in W$
    साठी $\mSat{M}{\psi}[w']$}
  \end{enumerate}
\end{defn}

इथे मात्र प्राप्यता संबंधांवर घालायच्या अटींचा विचार
करावा लागतो. खंड~\ref{aml:el:trw:defn:sub:mmodels-box} नुसार,
$R_a ww'$ असलेले एकही~$w'$ नसेल, तेव्हा
सूत्र~$\Knows_a \psi$ हे~$w$ येथे सत्य असते.
हा खंड सामान्य मोडल तर्कशास्त्रातील खंडासारखाच
आहे: एखाद्या जगाला उत्तरवर्ती जगे नसतील, तर सर्व
$\Box$-सूत्रे तेथे रिक्तपणे सत्य असतात. $\Knows$ हा
\emph{ज्ञान} दर्शवतो असे मानल्यास हे अत्यंत
प्रतिसहजभासी वाटते. कारण एखाद्या कर्त्याला एकाच
वेळी $\varphi$ आणि $\lnot \varphi$ ही दोन्ही माहीत असू शकतील,
अशी \emph{कोणतीही} परिस्थिती नसते, असे आपण
सहसा मानतो.

एक उपाय म्हणजे ज्ञानविषयक तर्कशास्त्रातील प्राप्यता
संबंध नेहमी \emph{परावर्ती} असेल याची खात्री करणे.
यातून वास्तविक जग कर्त्याच्या माहितीशी सुसंगत आहे,
ही कल्पना ढोबळमानाने व्यक्त होते. वस्तुतः ज्ञानविषयक
तर्कशास्त्रांत सहसा S5 वापरतात; पण $\Knows_a$ संबंधाने
नेमके काय दर्शवायचे आहे यानुसार काही जण त्याहून
दुर्बल व्यवस्था वापरू शकतात.

\begin{figure}
  \begin{center}
    \begin{tikzpicture}[modal]
      \node[world] (w1) [label={[align=right]right:\mTrue{p}\\ \mFalse{q}}]{$w_1$};
      \node[world] (w2) [label={[align=right]right:\mTrue{p}\\ \mTrue{q}},
        above right=of w1]{$w_2$};
      \node[world] (w3) [label={[align=right]right:\mFalse{p}\\ \mFalse{q}},
        below right=of w1] {$w_3$};
      \draw[<->] (w1) to node {$a$} (w2);
      \draw[->,loop left] (w1) to node {$a,b$} (w1);
      \draw[<->] (w1) to [swap] node {$b$} (w3);
      \draw[->,loop above] (w2) to node {$a,b$} (w2) ;
      \draw[->,loop below] (w3) to node {$a,b$} (w3) ;
    \end{tikzpicture}
  \end{center}
  \caption{साधे ज्ञानविषयक प्रतिमान.}
  \label{aml:el:trw:fig:simple}
\end{figure}

\begin{prob}
  \ref{aml:el:trw:fig:simple} मधील प्रतिमानात
  पुढीलपैकी कोणती विधाने सत्य आहेत ते ठरवा:
  \begin{enumerate}
    \item $\mSat{M}{\lnot q}[w_1]$;
    \item $\mSat{M}{\Knows_a \lnot q}[w_1]$;
    \item $\mSat{M}{\Knows_b \lnot q}[w_1]$;
    \item $\mSat{M}{\Knows_b q \lor \Knows_b \lnot q}[w_2]$;
    \item $\mSat{M}{\Knows_a( \Knows_b q \lor \Knows_b \lnot q)}[w_2]$;
    \item $\mSat{M}{\EKnows_{\{a,b\}} \lnot q}[w_3]$;
    \end{enumerate}
\end{prob}

जगात सत्य असण्याची मूलभूत व्याख्या दिल्यानंतर,
मोडल वैधता आणि तार्किक निष्पन्नता यांसारख्या सामान्य
मोडल तर्कशास्त्रातील इतर चिन्हार्थमीमांसक संकल्पनाही
इथे लागू होतात; फक्त मोडल संकारकांच्या अर्थलक्षणाकडे
पाहण्याचा हा नवा दृष्टिकोन वापरावा लागतो.

आता सामायिक-ज्ञान संकारक~$\CKnows_G$ साठी सत्यतेच्या
अटीही देता येतात. \ref{sfr:rel:ops:sec} मधील
व्याख्येनुसार, संबंध~$R$ चे \emph{संक्रमक संवरण}~$R^+$
पुढीलप्रमाणे असते:
\begin{align*}
  R^+ &= \bigcup_{0 < n \in \mathbb{N}} R^n,
\intertext{येथे}
  R^1 & = R \text{ आणि}\\
  R^{n+1} & = \Setabs{\tuple{x, z}}{\exists y (R^n xy \land Ryz)}\quad(n\geq1).
\end{align*}
मग $G$ हा कर्त्यांचा समूह असेल, तर सर्व कर्त्यांच्या
प्राप्यता संबंधांच्या संघाचे संक्रमक संवरण म्हणून
$R_G = ( \bigcup_{b
\in G} R_b )^+$ परिभाषित करतो.
यात एक किंवा अधिक प्राप्यता-पायऱ्या असलेले मार्ग घेतले
आहेत; शून्य पायऱ्यांचा मार्ग आपोआप जोडलेला नाही.

\begin{defn}
$G' \subseteq G$ असेल, तर $\mSat{M}{\CKnows_{G'} \varphi}[ w]$
तेव्हा आणि केवळ तेव्हाच, जेव्हा $R_{G'} w w'$ असलेल्या
प्रत्येक $w'$ साठी $\mSat{M}{\varphi}[w']$.
\end{defn}
ही अट सर्व रिकामेतर सीमित कर्तृक्रमांसाठी ज्ञानाची
अट घेण्यास समतुल्य आहे. येथे सामान्य $R_a$ संबंधांसाठी
संक्रमक संवरण घेतले आहे; परावर्ती संक्रमक संवरण नाही.
समूह रिकामेतर आणि त्यातील संबंध परावर्ती असतील, तर
$R_{G'}ww$ असल्याने सामायिक ज्ञान असलेले सूत्र~$w$
येथेही सत्य असते. रिकाम्या समूहासाठी संबंध रिकामा
आणि ही सार्विक अट रिक्तपणे सत्य असते.













\section{प्राप्यता संबंध आणि ज्ञानविषयक तत्त्वे}\label{aml:el:acc:sec}

सामान्य मोडल तर्कशास्त्रांतील चौकट-अनुरूपतेबद्दल आपल्याला
आधीच जे माहीत आहे, त्या आधारावर ज्ञानविषयक अर्थलक्षण
दिल्यास वैशिष्ट्यदर्शक सूत्रे कशी दिसतात हे
पाहता येईल. ज्ञानविषयक तर्कशास्त्रांचा अर्थ सहसा S5
मध्ये लावतात, हे आपण सांगितले आहे. आता विविध
ज्ञानविषयक तत्त्वे कशी व्यक्त होतात आणि त्यांचा
विविध चौकट-अटींशी कसा अनुरूप संबंध आहे ते पाहू.

सामान्य मोडल तर्कशास्त्रात विविध मोडल सूत्रे
प्राप्यता संबंधांचे निरनिराळे गुणधर्म वैशिष्ट्यित
करतात, हे आठवा. विशिष्ट ज्ञानविषयक तत्त्वांशी
संबंधित अशी काही सूत्रे या तक्त्यात दिली आहेत.

\begin{table}[t]
    \begin{tabular}{| p{.465\textwidth} || p{.465\textwidth} |}
      \hline
      {\emph{$R$ असा असेल \dots}} & {\emph{तर~$\mModel{M}$ मध्ये \dots सत्य असते:}} \\
      \hline \hline

      & $\Knows (p \lif q) \lif (\Knows p \lif \Knows q)$
      \hfill \newline{(संवरण)} \\
      \hline
      \emph{परावर्ती}: $\forall w Rww$
      & $\Knows p \lif p$ \hfill \newline{(ज्ञानाची सत्यता)} \\
      \hline
      \emph{संक्रमक}: \newline
      $\forall u \forall v \forall w ((Ruv \land Rvw) \lif Ruw)$ &
      $\Knows p \lif \Knows \Knows p$ \hfill
           \newline{(सकारात्मक आत्मनिरीक्षण)} \\
      \hline
      \emph{युक्लिडी}: \newline
      $\forall w \forall u \forall v ((Rwu \land Rwv) \lif Ruv)$ &
      $\lnot \Knows p \lif \Knows \neg \Knows p$ \hfill
        \newline{(नकारात्मक आत्मनिरीक्षण)} \\
      \hline
    \end{tabular}
    \caption{चार ज्ञानविषयक तत्त्वे.}
    \label{aml:el:acc:tab:four}
  \end{table}

$T$ स्वयंसिद्धाशी संबंधित ज्ञानाची सत्यता हे या
तत्त्वांपैकी सर्वात कमी विवाद्य मानतात, कारण
सूत्र माहीत असेल, तर ते सत्य असलेच पाहिजे,
असा तिचा आशय आहे. संवरण तसेच सकारात्मक आणि
नकारात्मक आत्मनिरीक्षण ही तत्त्वे त्यापेक्षा बरीच
अधिक विवाद्य आहेत.

$K$ स्वयंसिद्धाशी संबंधित संवरण म्हणजे कर्त्याचे
ज्ञान निहितार्थाखाली बंद असते, ही कल्पना. काही
प्रसंगी ती रास्त वाटू शकते. उदाहरणार्थ, मी
व्हिक्टोरियात असेन, तर व्हॅन्कुव्हर बेटावर आहे,
हे मला माहीत असू शकते. काही विलक्षण संशयवादी
परिस्थिती बाजूला ठेवल्या, तर मी व्हिक्टोरियात
आहे हे मला माहीत आहे; म्हणून मी व्हॅन्कुव्हर
बेटावर आहे हेही मला माहीत असावे. या बाबतीत
माझ्या ज्ञानाचे तार्किक संवरण बऱ्यापैकी सहज
पटण्यासारखे वाटेल. दुसरीकडे, आपल्या ज्ञानाचे
परिणाम आपण नेहमी विचारात घेत नाही; म्हणून इतर
प्रसंगी हा निष्कर्ष तितका सहज पटणार नाही.

सकारात्मक आत्मनिरीक्षणाला कधीकधी KK-तत्त्व म्हणतात:
मला एखादी गोष्ट माहीत असेल, तर ती मला माहीत आहे
हेही मला माहीत असते, असे ते सांगते. ते 4 स्वयंसिद्धाचे
ज्ञानविषयक समरूप आहे. त्याचप्रमाणे नकारात्मक
आत्मनिरीक्षण असे सांगते की एखादी गोष्ट मला
\emph{माहीत नसेल}, तर ती मला माहीत नाही हे मला
माहीत असते; ते 5 स्वयंसिद्धाचे समरूप आहे. या
तत्त्वांतील फरक प्रतिउदाहरणातही जपावा. सकारात्मक
आत्मनिरीक्षण मोडण्यासाठी एखादी गोष्ट मला माहीत असूनही
ती मला माहीत आहे हे माहीत नसणे आवश्यक आहे.
नकारात्मक आत्मनिरीक्षण मोडण्यासाठी ती गोष्ट मला
माहीत नसूनही, ती मला माहीत नाही हे माहीत नसणे
आवश्यक आहे. प्रत्यक्षात माहीत असलेल्या गोष्टीबद्दल
साशंक असणे हे पहिल्या तत्त्वासाठी संबंधित ठरते;
दुसऱ्याच्या पूर्वांगात मात्र माहीत नसणे येते.














\section{द्विअनुकरणे}\label{aml:el:bsd:sec}

आपल्या ज्ञानविषयक भाषेच्या अभिव्यक्तिक्षमतेविषयी
उरलेला एक प्रश्न प्रतिमाने आणि त्यांत सत्य असलेली
सूत्रे यांमधील संबंधाशी निगडित आहे. चौकट-अनुरूपतेच्या
निष्कर्षांवरून आपण पाहिले की काही सूत्रे एखाद्या
चौकटीत वैध असतील, तर त्या चौकटीला विशिष्ट गुणधर्म
असतात. पण उदाहरणार्थ, एखाद्या जगापासून त्याच
जगाकडे जाणारा परावर्ती बाण आणि प्रत्येक जग पुढच्या
जगाकडे जाणारी अनंत साखळी यांमध्ये भेद करता येईल
अशी आपली मोडल भाषा आहे का? म्हणजे, या दोन
जगांपैकी फक्त एकाच जगात सत्य असेल असे एखादे
सूत्र~$A$ आहे का?

संबंधी प्रतिमानांतील संबंधित जगांत मूलभूत
ज्ञानविषयक सूत्रांनी भेद करता येत नाही, ही कल्पना
द्विअनुकरण संबंध पकडतो. यासाठी दोन प्रतिमानांतील
जगांची संख्या किंवा त्यांचे ग्राफ समरूप असण्याची
अट नाही; पुढील अटी योग्य अनुरूपता सांगतात.

\begin{defn}[द्विअनुकरण]
$M_1 = \tuple{W_1, R_1, V_1}$ आणि $M_2 = \tuple{W_2, R_2, V_2}$
ही एकाच कर्तृचिन्ह-संच~$G$ व एकाच विधानीय
भाषेची दोन संबंधी प्रतिमाने असू द्या. तसेच
$\mathcal{R} \subseteq W_1 \times W_2$ हा द्विस्थानी
संबंध असू द्या. प्रत्येक $\tuple{w_1, w_2} \in
\mathcal{R}$ साठी पुढील अटी पूर्ण होत असतील, तर
$\mathcal{R}$ हे \emph{द्विअनुकरण} आहे असे म्हणतो:

\begin{enumerate}
  \item सर्व विधानीय चले~$p$ साठी
  $w_1 \in V_1(p)$ तेव्हा आणि केवळ तेव्हाच, जेव्हा
  $w_2 \in V_2(p)$.
  \item सर्व कर्ते $a \in G$ आणि सर्व जगे $v_1 \in W_1$
  यांसाठी, $R_{1_a} w_1 v_1$ असेल, तर $R_{2_a} w_2 v_2$
  आणि $\tuple{v_1, v_2} \in \mathcal{R}$ असलेले काही
  $v_2 \in W_2$ असते.
  \item सर्व कर्ते $a \in G$ आणि सर्व जगे $v_2 \in W_2$
  यांसाठी, $R_{2_a} w_2 v_2$ असेल, तर $R_{1_a} w_1 v_1$
  आणि $\tuple{v_1, v_2} \in \mathcal{R}$ असलेले काही
  $v_1 \in W_1$ असते.
\end{enumerate}

$M_1$ आणि $M_2$ यांमधील एखादे द्विअनुकरण जगे
$w_1$ आणि $w_2$ यांना जोडत असेल, तर आपण
$\tuple{M_1, w_1} \leftrightarroweq
\tuple{M_2, w_2}$ असेही लिहू शकतो. तेव्हा
$\tuple{M_1, w_1}$ आणि $\tuple{M_2, w_2}$ यांना
\emph{द्विअनुकरणसदृश} म्हणतो.
\end{defn}

द्विअनुकरण संबंधातील निरनिराळ्या अटी वेगवेगळ्या
गोष्टी सुनिश्चित करतात. पहिली अट सर्व विधानीय चले
बाबतीत एकमत असल्याचे सुनिश्चित करते;
म्हणून द्विअनुकरणसदृश जगांत तीच मोडल-विरहित सूत्रे
सत्य असतात. उरलेल्या दोन अटींना अनुक्रमे ``पुढे''
आणि ``मागे'' अशा अटी म्हणतात. संबंधित जगांच्या प्रत्येक
प्राप्य उत्तरवर्ती जगाला दुसऱ्या प्रतिमानात संबंधित
प्राप्य उत्तरवर्ती जग दोन्ही दिशांनी मिळते; उत्तरवर्ती
जगांची संख्या सारखी किंवा बाणांच्या रचना समरूप असण्याची अट नाही.

\begin{thm}
$\tuple{M_1, w_1} \leftrightarroweq \tuple{M_2, w_2}$
असेल, तर प्रत्येक सूत्र~$\varphi$ साठी
$\mSat{M_1}{\varphi}[w_1]$ तेव्हा आणि केवळ तेव्हाच,
जेव्हा $\mSat{M_2}{\varphi}[w_2]$.
\end{thm}

\begin{figure}
  \begin{center}
    \begin{tikzpicture}[modal]
      \node[world] (w1) {$w_1$};
      \node[world] (w2) [above left=of w1]{$w_2$};
      \node[world] (w3) [above right=of w1] {$w_3$};
      \draw[<->] (w1) to node {$a$} (w2);
      \draw[->,loop below] (w1) to node {$a$} (w1);
      \draw[<->] (w1) to [swap] node {$a$} (w3);
      \draw[->,loop above] (w2) to node {$a$} (w2) ;
      \draw[->,loop above] (w3) to node {$a$} (w3) ;

      \node[world](v1)[right of=w1, xshift=1.5in]{$v_1$};
      \node[world](v2)[above of=v1]{$v_2$};
      \draw[<->] (v1) to node {$a$} (v2);
      \draw[->,loop below] (v1) to node {$a$} (v1);
      \draw[->,loop above] (v2) to node {$a$} (v2) ;

      \draw[-,dotted] (w1) to (v1) ;
      \draw[-,dotted,bend left=45] (w2) to (v2) ;
      \draw[-,dotted,bend left=30] (w3) to (v2) ;
    \end{tikzpicture}
  \end{center}
  \caption{दोन द्विअनुकरणसदृश प्रतिमाने.}
  \label{aml:el:bsd:fig:bisimilar}
\end{figure}

\ref{aml:el:bsd:fig:bisimilar} मध्ये बाणांची रचना दाखवली आहे;
द्विअनुकरणासाठी मूल्यांकनांची अनुरूपताही आवश्यक आहे.
येथे प्रत्येक विधानीय चल~$p$ साठी
$w_1$ आणि $v_1$ येथे त्याचे सत्यतामूल्य सारखे,
आणि $w_2,w_3,v_2$ या तिन्ही जगांत त्याचे सत्यतामूल्य
सारखे घेऊ. इतर कर्ते असतील, तर त्यांचे संबंध या
दोन प्रतिमानांत रिकामे घेऊ. मग
$\{\tuple{w_1,v_1},\tuple{w_2,v_2},\tuple{w_3,v_2}\}$
हा संबंध द्विअनुकरण आहे. म्हणून संबंधित जगांत मूलभूत
मोडल भाषेतील सूत्रांनी भेद करता येत नाही. मात्र
$w_2$ आणि~$w_3$ यांत निरनिराळी विधानीय चले
सत्य असतील, तर ते दोन्ही~$v_2$ शी जोडणारा हाच
संबंध पहिली अट पूर्ण करत नाही.












\section{सार्वजनिक घोषणा तर्कशास्त्र}\label{aml:el:pal:sec}

गतिशील ज्ञानविषयक तर्कशास्त्रांमुळे कर्त्यांचे ज्ञान
काळानुसार किंवा नवी माहिती मिळाल्यावर कसे बदलते हे
दर्शवता येते. यांपैकी अनेक तर्कशास्त्रे ज्ञानातील बदल
माहितीपर \emph{घटना} किंवा \emph{अद्यतने} यांद्वारे
दर्शवतात. अद्यतनाचा सर्वांत मूलभूत प्रकार म्हणजे
सार्वजनिक घोषणा: एखादे सूत्र सत्य असल्याची घोषणा
होते आणि सर्व कर्ते ती एकत्रितपणे पाहतात. यासाठी
भाषेचा पुढीलप्रमाणे विस्तार करतो.

\begin{defn}
$G$ हा कर्तृचिन्हांचा संच असू द्या. सार्वजनिक घोषणा
असलेल्या अनेक कर्त्यांच्या ज्ञानविषयक तर्कशास्त्राच्या
मूलभूत भाषेत पुढील गोष्टी असतात:
\begin{enumerate}
  \item असत्यता दर्शवणारा विधानीय स्थिरांक~$\lfalse$.
  \item सत्यता दर्शवणारा विधानीय स्थिरांक~$\ltrue$.
  \item विधानीय चलांचा गणनीय अनंत संच: $\Obj
    p_0$, $\Obj p_1$, $\Obj p_2$, \dots
  \item विधानीय संयोजके: \startycommalist
  \ycomma $\lnot$ (नकार)
  \ycomma $\land$ (संयोग)
  \ycomma $\lor$ (विकल्पयोग)
  \ycomma $\lif$ (शर्त)
  \ycomma $\liff$ (द्विशर्त).
  \item $a \in G$ असताना ज्ञान-संकारक $\Knows_a$.
  \item $\psi$ हे सूत्र असताना सार्वजनिक घोषणा-संकारक $[\psi]$.
\end{enumerate}
\end{defn}

सार्वजनिक घोषणा-संकारक बॉक्स-संकारकासारखे कार्य
करतो. त्यानुसार भाषेची आगमनाधारित व्याख्या पुढीलप्रमाणे:

\begin{defn}
  ज्ञानविषयक भाषेतील \emph{सूत्रे} ची पुढीलप्रमाणे
  आगमनाने व्याख्या करतो:
  \begin{enumerate}
  \item $\lfalse$ हे अण्विक सूत्र आहे.

  \item $\ltrue$ हे अण्विक सूत्र आहे.

  \item प्रत्येक विधानीय चल~$\Obj p_i$ हे (अण्विक)
  सूत्र आहे.

  \item $\varphi$ हे सूत्र असेल, तर $\lnot \varphi$
  हेही सूत्र आहे.

  \item $\varphi$ आणि $\psi$ ही सूत्रे असतील, तर $(\varphi \land
  \psi)$ हे सूत्र आहे.

  \item $\varphi$ आणि $\psi$ ही सूत्रे असतील, तर $(\varphi \lor \psi)$
  हे सूत्र आहे.

  \item $\varphi$ आणि $\psi$ ही सूत्रे असतील, तर $(\varphi \lif \psi)$
  हे सूत्र आहे.

  \item $\varphi$ आणि $\psi$ ही सूत्रे असतील, तर $(\varphi \liff \psi)$
  हे सूत्र आहे.

  \item $\varphi$ हे सूत्र असेल आणि $a \in G$ असेल,
  तर $\Knows_a \varphi$ हे सूत्र आहे.

  \item $\varphi$ आणि $\psi$ ही सूत्रे असतील,
  तर $[\varphi] \psi$ हे सूत्र आहे.

  \item याखेरीज दुसरे काहीही सूत्र नाही.
\end{enumerate}
\end{defn}

सूत्र $[\varphi] \psi$ चा अभिप्रेत अर्थ ``$\varphi$ सत्य
असल्याची सार्वजनिक घोषणा झाल्यानंतर $\psi$~सत्य असते''
असा आहे. सार्वजनिक घोषणांच्या संदर्भात सामायिक
ज्ञानाविषयी बोलणेही कधी उपयोगी ठरते; म्हणून भाषेत
सामायिक-ज्ञान संकारक~$\CKnows_G
\varphi$ देखील समाविष्ट करता येतो.













\section{सार्वजनिक घोषणा तर्कशास्त्राचे अर्थनिर्धारण}\label{aml:el:psm:sec}

सार्वजनिक घोषणा तर्कशास्त्रांची संबंधी प्रतिमाने
ज्ञानविषयक तर्कशास्त्रांतील प्रतिमानांसारखीच असतात.
मात्र सार्वजनिक घोषणा-संकारकाचे अर्थनिर्धारण नवे आहे.

\begin{defn}\label{aml:el:psm:defn:mmodels} प्रतिमान~$\mModel M = \tuple{W, R, V}$
  मध्ये \emph{सूत्र~$\varphi$ हे~$w$ येथे सत्य असणे},
  म्हणजे $\mSat{M}{\varphi}[w]$, पुढीलप्रमाणे आगमनाने
  परिभाषित करतो:
  \begin{enumerate}
  \item \indcase{\varphi}{\lfalse}{$\mSat{M}{\lfalse}[w]$ कधीच नसते}.
  \item \indcase{\varphi}{\ltrue}{$\mSat{M}{\ltrue}[w]$ नेहमी असते}.
  \item $\mSat{M}{p}[w]$ तेव्हा आणि केवळ तेव्हाच, जेव्हा $w \in V(p)$
  \item \indcase{\varphi}{\lnot \psi}{$\mSat{M}{\indfrm}[w]$
    तेव्हा आणि केवळ तेव्हाच, जेव्हा $\mSat/{M}{\psi}[w]$}.
  \item \indcase{\varphi}{(\psi \land \chi)}{$\mSat{M}{\indfrm}[w]$
    तेव्हा आणि केवळ तेव्हाच, जेव्हा $\mSat{M}{\psi}[w]$ आणि
    $\mSat{M}{\chi}[w]$}.
  \item \indcase{\varphi}{(\psi \lor \chi)}{$\mSat{M}{\indfrm}[w]$
    तेव्हा आणि केवळ तेव्हाच, जेव्हा $\mSat{M}{\psi}[w]$ किंवा
    $\mSat{M}{\chi}[w]$} (किंवा दोन्ही).
  \item \indcase{\varphi}{(\psi \lif \chi)}{$\mSat{M}{\indfrm}[w]$
    तेव्हा आणि केवळ तेव्हाच, जेव्हा $\mSat/{M}{\psi}[w]$ किंवा
    $\mSat{M}{\chi}[w]$}.
  \item \indcase{\varphi}{(\psi \liff \chi)}{$\mSat{M}{\indfrm}[w]$
    तेव्हा आणि केवळ तेव्हाच, जेव्हा $\mSat{M}{\psi}[w]$ आणि
    $\mSat{M}{\chi}[w]$ ही दोन्ही असतात, किंवा
    $\mSat{M}{\psi}[w]$ आणि $\mSat{M}{\chi}[w]$ यांपैकी
    कोणतेही नसते}.
  \item\label{aml:el:psm:defn:sub:mmodels-box}
    \indcase{\varphi}{\Knows_a \psi}{$\mSat{M}{\indfrm}[w]$ तेव्हा आणि
    केवळ तेव्हाच, जेव्हा $R_a ww'$ असलेल्या प्रत्येक $w' \in W$
    साठी $\mSat{M}{\psi}[w']$}
  \item\label{aml:el:psm:defn:sub:mmodels-pal}
    \indcase{\varphi}{[\psi] \chi}{$\mSat{M}{\indfrm}[w]$ तेव्हा आणि
    केवळ तेव्हाच, जेव्हा एकतर $\mSat/{M}{\psi}[w]$,
    किंवा $\mSat{M}{\psi}[w]$ आणि $\mSat{M \mid \psi}{\chi}[w]$}


  येथे $\mModel M \mid \psi = \tuple{W', R', V'}$ ची
  व्याख्या पुढीलप्रमाणे आहे:
  \begin{enumerate}
    \item $W' = \Setabs{ u \in W}{\mSat{M}{\psi}[u]}$.
      म्हणजे $\mModel M \mid \psi$ मधील जगे ही~$\mModel M$
      मध्ये $\psi$ सत्य असलेली जगेच आहेत.
    \item $R'_a = R_a \cap (W' \times W')$. प्रत्येक कर्त्याचा
      प्राप्यता संबंध केवळ~$W'$ मध्ये उरलेल्या जगांपुरता
      मर्यादित केला जातो.
    \item $V'(p) = \Setabs{ u \in W'}{u \in V(p) }$.
      त्याचप्रमाणे उरलेल्या जगांतील विधानीय मूल्यांकने
      तीच राहतात. माहितीपर घटनांमुळे विधानीय चलांचे
      सत्यतामूल्य बदलत नाही, ही कल्पना त्यातून दिसते.
  \end{enumerate}
  $W'$ रिकामा असू शकतो. पण $\psi$ हे~$w$ येथे असत्य
  असेल, तर घोषणा-सूत्र पहिल्या पर्यायाने सत्य आहे आणि
  त्या प्रतिबंधित रचनेत~$w$ येथे $\chi$ तपासायचे नाही.
  $\psi$ हे~$w$ येथे सत्य असेल, तर $w\in W'$ असल्याने
  $W'$ रिक्तेतर आहे आणि प्रतिबंधित प्रतिमानातील
  पुढील मूल्यांकन नीट परिभाषित आहे.

  \end{enumerate}
\end{defn}

म्हणून सार्वजनिक घोषणा तर्कशास्त्रांचे वैशिष्ट्य असे
की, प्रतिमान~$\mModel M$ येथे सूत्र सत्य आहे
की नाही हे कधीकधी~$\mModel M$ व्यतिरिक्त दुसरे
प्रतिमान पाहिल्याशिवाय ठरवता येत नाही.

हेही लक्षात घ्या की आपल्या अर्थनिर्धारणात घोषणा-संकारक
हा $\Box$~संकारक आहे. म्हणून एखाद्या जगात
सूत्र~$\varphi$ सत्य असल्याची घोषणा करता येत
नसेल, तर $[\varphi] \psi$ तेथे आपोआप सत्य असते; जसे
अंतिम जगांत सर्व $\Box$ सूत्रे सत्य असतात.

\begin{figure}
  \begin{center}
    \begin{tikzpicture}[modal]
      \node[world] (w1) [label={below left:$p, \lnot q$}]{$w_1$};
      \node[world] (w2) [left=of w1, label={left:$\lnot p, \lnot q$}]{$w_2$};
      \node[world] (w3) [above=of w1, label={right:$p, q$}] {$w_3$};
      \draw[<->] (w1) to [swap] node {$b$} (w2);
      \draw[->,loop below] (w1) to node {$a, b$} (w1);
      \draw[<->] (w1) to [swap] node {$a$} (w3);
      \draw[->,loop above] (w2) to node {$a, b$} (w2) ;
      \draw[->,loop above] (w3) to node {$a, b$} (w3) ;

      \node[world](v1)[right of=w1, xshift=1.25in, label={right:$p, \lnot q$}]{$w'_1$};
      \node[world](v3)[above of=v1,label={right:$p,q$}]{$w'_3$};
      \draw[<->] (v1) to node {$a$} (v3);
      \draw[->,loop below] (v1) to node {$a,b$} (v1);
      \draw[->,loop above] (v3) to node {$a,b$} (v3) ;

      \node(m)[below=of w1]{$\mModel M$};
      \node(m')[below=of v1]{$\mModel M \mid p$};

      \draw[-,dotted] (w1) to node {\footnotesize $p$ ची घोषणा}(v1) ;

    \end{tikzpicture}
  \end{center}
  \caption{$p$ ची सार्वजनिक घोषणा होण्यापूर्वी आणि झाल्यानंतर.}
  \label{aml:el:psm:fig:announcement-example}
\end{figure}

सूत्राची सार्वजनिक घोषणा म्हणजे प्रतिमान
संकुचित करणे, किंवा ज्या जगांत ते सूत्र सत्य
होते तितक्याच जगांपुरते ते मर्यादित करणे, असे पाहता
येते. \ref{aml:el:psm:fig:announcement-example} मध्ये~$p$ ची
सार्वजनिक घोषणा केल्याचा परिणाम दाखवला आहे.
या प्रतिमानाचे एक वैशिष्ट्य असे की घोषणेमुळे
कर्ता~$b$ ला~$p$ माहीत होते, पण कर्ता~$a$ ला
तसे नव्याने कळत नाही (कारण~$p$ सत्य आहे हे
$a$ ला आधीच माहीत होते).

अधिक औपचारिकपणे, $\mSat{M}{\lnot \Knows_b p}[w_1]$
असते; पण $\mSat{M
\mid p}{\Knows_b p}[w'_1]$ असते. यावरून
$\mSat{M}{[p] \Knows_b
p}[w_1]$ मिळते. पण घोषणेच्या परिणामाबद्दल
याहूनही अधिक प्रबळ विधान करता येते. प्रत्यक्षात
$\mSat{M}{[p]\CKnows_{\{a,b\}} p}[w_1]$ असते.
म्हणजे $p$ ची घोषणा झाल्यावर ते \emph{सामायिक ज्ञान}
बनते.

पण हे सर्वसाधारणपणे खरे आहे का? $\varphi$ सत्य असल्याची
घोषणा केल्यावर $\varphi$ हे \emph{नेहमीच} सामायिक ज्ञान बनेल
का? आश्चर्य वाटेल, पण उत्तर नाही असे आहे. प्रत्यक्षात,
काही सूत्रे सत्य असताना त्यांची घोषणा करणे
शक्य आहे, आणि घोषणेनंतर ती सत्य राहत नाहीत.
उदाहरणार्थ, \ref{aml:el:psm:fig:announcement-example} मधील
$w_1$ येथे $p \land \lnot \Knows_b p$ ची घोषणा
केल्याचा परिणाम पाहा. येथे $p$ हे $w_1,w_3$ या
दोन्ही जगांत सत्य आहे, म्हणून $\mModel M\mid p$ मध्ये
ही दोन्ही जगे उरतात. पण $w_3$ येथे $b$ चा एकमेव
प्राप्य पर्याय $w_3$ असल्याने $\Knows_b p$ सत्य आहे;
त्यामुळे $p\land\lnot\Knows_b p$ तेथे असत्य आहे.
हे संपूर्ण सूत्र फक्त $w_1$ येथे सत्य असते. म्हणून
$\mModel M\mid(p\land\lnot\Knows_b p)$ मध्ये केवळ
$w_1$ उरते आणि दोन्ही कर्त्यांसाठी फक्त त्याचे स्वबाण
उरतात; ते $\mModel M\mid p$ सारखे प्रतिमान नाही.
या एक-जगी प्रतिमानात $p$ आणि $\Knows_b p$ दोन्ही
सत्य आहेत. म्हणून
$\mSat{M \mid (p \land \lnot \Knows_b p)}{\lnot
(p \land \lnot \Knows_b p)}[w_1]$ असते. म्हणजे
घोषणा केल्यावर असत्य होणारे हे सूत्र आहे.



\part{अंतःप्रज्ञावादी तर्कशास्त्र}\label{int:part}
